\scanpage{55}
\mspage{7}
avec le formulaire \unsure{familier|géométrique} dans $\widetilde\Pi_{0,3}\simeq\unsure{\operatorname{SL}(2,\mathbb Z)|\operatorname{GL}(2,\mathbb Z)}$ :
\[
(33)\quad\left\{\begin{aligned}
\varepsilon_0&\overset{\mathrm{def}}=\sigma_\infty\rho\\
\varepsilon_1&\overset{\mathrm{def}}=\rho(\varepsilon_0)\quad(=\rho\varepsilon_0\rho^{-1})\\
\varepsilon_\infty&\overset{\mathrm{def}}=\rho(\varepsilon_1)=\rho^2(\varepsilon_0)=\rho^{-1}(\varepsilon_0)
\end{aligned}\right.
\qquad
\begin{aligned}
l_0&\overset{\mathrm{def}}=\varepsilon_0^2\\
l_1&\overset{\mathrm{def}}=\varepsilon_1^2\\
l_\infty&\overset{\mathrm{def}}=\varepsilon_\infty^2
\end{aligned}
\]
\[
l_\infty l_1l_0=1
\]
\margin{
$\rho(l_i)=l_{\rho(i)}\quad\left\{\begin{aligned}\rho(l_0)&=l_1\\\rho(l_1)&=l_\infty\\\rho(l_\infty)&=l_0\end{aligned}\right.$\newline
$\rho(\unsure{\varepsilon_i|l_i'})=\varepsilon_{\rho(i)}$\newline
\uline{NB} on fait opérer $\rho$ sur $0,1,\infty$ par $0\mapsto1\mapsto\infty\mapsto0$
}
\[
(33)\quad\left\{\begin{aligned}
\tau_\infty&\overset{\mathrm{def}}=\sigma_\infty\widetilde\sigma_\infty\unsure{=|,}\widetilde\sigma_\infty\sigma_\infty\\
\tau_\infty^2&=1
\end{aligned}\right.
\]
\[
\begin{aligned}
\tau_\infty(l_0)&=l_0^{-1}\\
\tau_\infty(l_1)&=l_1^{-1}\\
\tau_\infty(l_\infty)&=\unsure{l_0|l_0^{-1}}l_1=\unsure{\operatorname{int}(l_0)}(l_\infty^{-1})
\end{aligned}
\qquad
\begin{aligned}
\widetilde\sigma_\infty(l_0)&=l_1^{-1}\\
\widetilde\sigma_\infty(l_1)&=l_0^{-1}\\
\widetilde\sigma_\infty(l_\infty)&=l_\infty^{-1}
\end{aligned}
\]
\[
\begin{aligned}
\sigma_\infty(l_0)&=l_1\\
\sigma_\infty(l_1)&=l_0\\
\sigma_\infty(l_\infty)&=l_1^{-1}l_0^{-1}=\unsure{\operatorname{int}(l_0)}(l_\infty)
\end{aligned}
\]
Notons qu'un épinglage de la situation est donné par un élément du produit
\[
\operatorname{Rep}(J,\uline{\unsure{\bar\omega|\varpi}})\simeq J\times\uline{\omega}\times\uline{\unsure{\bar\omega|\varpi}}
\]
donc par la donnée d'un $(i,\omega,\unsure{\bar\omega|\varpi})$, définissant \deleted{donc} un $\varepsilon=\omega\wedge\unsure{\bar\omega|\varpi}\in\uline{\varepsilon}\simeq\uline{\omega}\wedge\uline{\unsure{\bar\omega|\varpi}}$. On va définir pour
\[
r=(i,\omega,\unsure{\bar\omega|\varpi})
\]
des isoms
\[
\deleted{\unsure{\Pi_{03}}}\qquad(34)\quad\left\{\begin{aligned}
\psi_{\deleted{\unsure{i}}\varepsilon}^{\deleted{\unsure{\omega}}i,\omega}&:\Pi_{03}\xrightarrow{\sim}\Pi_\varepsilon\\
\psi_{i\deleted{\unsure{\omega}}}^{\omega,\varepsilon}&:\Pi_{03}\xrightarrow{\sim}\Pi_i
\end{aligned}\right.
\]
\margin{isoms dits \uline{d'épinglage}}
donnant lieu comme de juste au triangle commutatif
\diagram{
\begin{tikzcd}[ampersand replacement=\&]
\& \Pi_{03}\arrow[dl,"\psi_{\deleted{\unsure{i}}\varepsilon}^{\deleted{\unsure{\omega}}i,\omega}"']\arrow[dr,"\psi_{i\deleted{\unsure{\omega}}}^{\omega,\varepsilon}"] \& \\
(35)\quad\Pi_\varepsilon\arrow[rr,shift left=.6ex,"\varphi_{i\varepsilon}","\sim"'] \& \& \Pi_i\arrow[ll,shift left=.6ex,"\varphi_{\varepsilon i}","\sim"']
\end{tikzcd}
}
\margin{54}
